\section{讨论}
\subsection{已经连接的部分与仍需区分的对象}
降阶模型连接了三维几何、变惯性、磁能、物理电流端口与隐式积分。在光滑全线圈反馈分支上，同一原向量场也用于输运形式和普通中点的比较。这是力学与局部结构之间具体可计算的连接，但精度判断必须分解：对截断形式精确匹配的坐标图，未必精确积分原场。所报告全线圈窗口很短，图构造成本也远大于所模拟的物理时长。

受限采样控制器及环境空间杆是进一步但不同的计算。其约束可行性与功率记录不能继承理想保守隐式映射的 exact-symplectic 证明。电流记忆、激活约束、插值与环境载荷改变了数学对象。同样，所学网络逼近的是理想偏置能量，而非完整线圈 Hamiltonian。学习到执行的已实现诊断是静态力分配；本文尚无将完整学习控制器经局部 Birkhoff 图积分、再把这一组合控制器转移到空间杆的实验。不能将各分支分别成功的结果相加，代替缺失的组合验证。

\subsection{物理及数学证据的局限}
仿射场的力偏差对应本文测试的几何与状态。它说明需要检验实际采用的近似，而不是给出所有仿射磁模型的普遍误差下界。完整线圈结果本身仍依赖磁材料假设并具有空间求积误差。12个电流通道并未消除测得的广义力秩亏。短时ROM控制比较能够分辨被控对象模型差异，却未能分辨控制器变化带来的跟踪收益。较长空间杆测试的跟踪略差，且观察到的时间细化阶较低。两者均不足以支持导航性能改善的主张。

实际应用目标仍广于已验证的计算层。ROM输运计算采用明确的光滑载荷，不含脉动血流与接触切换。空间杆采用规定血流和刚性管壁，以及合成残余加速度。双向流固耦合、柔性血管、生理边界数据、硬件线圈标定及微重力验证尚未完成。短时保守ROM--杆比较不能约束受载血管中的降阶误差，而且违反完整杆上的一项梯度限制。保留该诊断有助于避免把较小状态差异误读为可行介入。

数学结论也具有所述局部性或条件性。采样缺陷不能给出坐标图、逆映射、导数与扰动的统一界。冻结复网络缺少已验证的解析域误差预算，能量曲率失败限制了其机械解释。共振四次项计算提供第一频率对的候选结构，但尚未建立完整不变子系统或12作用量 twist 证书。本文未将连续或离散 KAM 定理应用于血流、接触、采样受限反馈或所学三维模型。

\section{结论}
三维几何杆能够在不冻结质量矩阵的情况下接入可计算的局部 Birkhoff 表示，也能够在不把任意相空间修正视作执行输入的情况下接入真实磁端口。已有结果量化了这些连接及其限制。高阶输运降低表示误差，保守 Central4 达到预期阶数，全线圈虚功与磁能求导一致。这些成立结果同时伴随显著仿射力误差、秩受限控制、学习动量能量非凸、尚未解决的共振持久性以及未得到空间杆跟踪改善。证据支持局部保结构仿真和物理一致的磁执行接口，但尚未建立统一的学习保结构血管控制器，也未给出三维 KAM 证书。

\section*{声明}
\paragraph{作者与贡献。} 作者姓名、单位、通讯作者信息及贡献声明须由负责任的人类作者在投稿前补充并批准。
\paragraph{基金。} 尚未提供资助来源和项目编号，须由作者补充。
\paragraph{利益冲突。} 尚未提供声明，须由作者在投稿前完成。
\paragraph{伦理与知情同意。} 本研究报告数学模型及数值仿真，不分析人类受试者、患者或动物数据。机构要求及最终声明仍由作者负责确认。
\paragraph{数据与代码可用性。} 来源记录位于 \url{https://github.com/mop74720-commits/1028} 的 Stage~5 提交 \texttt{1c73a58}，完整标识见随附 README。论文包提供源字段及 SHA-256 清单、绘图程序、冻结数据路径和双语补充证明。重建图表不会重新运行科学实验。缺失历史材料及被否定输出均予以明确标识，不以重新生成文件代替。
\paragraph{AI辅助准备披露。} OpenAI Codex 辅助进行了来源检查、论文组织、起草、翻译、数据提取、图表准备及 LaTeX 编译。其不作为作者。负责任的作者须核验数学、证据、文献与声明，并在投稿前批准两个语言版本。本次论文整理未新增科学实验。
